\newsec{Introduzione}{\vspace{-\bigskipamount}\vspace{-\bigskipamount}}

\subsection{Raggi X}

I raggi X sono una parte dello spettro della radiazione elettromagnetica: li distinguiamo dalle lunghezze d'onda comprese nell'intervallo $ 1 \unit{pm} - 1 \unit{\angstrom} $.

Il tubo radiogeno (di Röntgen) è uno dei primi strumenti utilizzato per produrre dei raggi X.
Un flusso di corrente su un filamento resistivo di tungsteno dissipa del calore: al crescere dell'intensità di corrente alcuni elettroni vengono estratti dal filamento (effetto termoionico).
Gli elettroni vengono accelerati da una grande differenza di potenziale $ V_\text{EHT} $, e diretti verso del rame.

I raggi X vengono generati per bremsstrahlung e per interazione con le shell di elettroni del rame (raggi caratteristici).

\subsection{Setup sperimentale}

I raggi X prodotti dal tubo Röntgen vengono collimati e diretti verso un contatore Geiger-Müller.
Questo detector contiene un gas a bassa pressione da cui la radiazione incidente ionizza/estrae (per effetto fotoelettrico) elettroni: questi generano dei procedimenti a valanga producendo sempre più elettroni.
Proprio per questo lo strumento è solo in grado di contare il numero di eventi di interazione, e non di misurare la loro energia (per esempio).
